\scanpage{32}
\uline{Th} Soit $G$ un préschéma en groupes \inserted{lisse} séparé de présentation finie sur $S$. Alors

\noindent\begin{tabular}{@{}|p{.94\linewidth}@{}}
(i) \deleted{les fonctions} \deleted{\unsure{Th|les}} les fonctions \inserted{$\rho_{aff},\deleted{\unsure{\rho}}d\deleted{\dim G}$ sont $1/2$ continues sup, et de} $\rho_n,\rho_{na},\rho_u$ sont semi-continues \uline{supérieurement} \deleted{si $G$ est plat sur $S$ [\unsure{alors|donc} $d$ est \unsure{lct.|loc.}\ \unsure{cte|c\textsuperscript{st}}]}.
\end{tabular}

\noindent\begin{tabular}{@{}|p{.94\linewidth}@{}}
(ii) \deleted{Si $G$ est \uline{plat} sur $S$,} \deleted{les} fonctions $\rho_{ab},\rho_{ab}+\rho_r,\rho_s$ sont semi-continues inférieurement.
\end{tabular}

On se ramène comme d'habitude au cas $S$ affine, \unsure{noth|noeth}, spectre d'un anneau de valuation discrète, pt fermé $s_0$ pt générique $s_1$, et à prouver les inégalités
\[
\left\{\begin{aligned}
&\inserted{\rho_{aff}\geq\rho'_{aff}\qquad d\geq d'}\\
&\deleted{\text{\textcircled{\scriptsize\unsure{4}}}}\text{\textcircled{\scriptsize 5}}\rho_{\deleted{\unsure{s}}n}\geq\rho'_n,\quad
\deleted{\text{\textcircled{\scriptsize\unsure{5}}}}\text{\textcircled{\scriptsize 6}}\rho_u\geq\rho'_u\quad
\text{\textcircled{\scriptsize 7}}\rho_{na}\geq\rho'_{na}\quad\deleted{,\rho_{aff}\geq\rho'_{aff},}\\
&\text{\textcircled{\scriptsize 1}}\rho_{ab}\leq\rho'_{ab},\quad \text{\textcircled{\scriptsize 2}}\rho_{ab}+\rho_r\leq\rho'_{ab}+\rho'_r,\quad \text{\textcircled{\scriptsize 3}}\rho_s\leq\rho'_s
\end{aligned}\right.
\]
\deleted{\unsure{les}}\unsure{les} \deleted{deux}\inserted{quatre} premières sans réserve de platitude.

Ce sont \deleted{\unsure{des}} (pour un groupe algébrique sur un corps alg clos $k$) des quantités invariantes en changeant $G$ en $G_{\mathrm{réd}}^\circ$, et de plus pour $\rho_{aff}$ et $d$, c'est une fonction \uline{\unsure{croissante|croissant}} du groupe. Donc cela nous permet [quitte à remplacer $G$ par $\overline{G_1^{\unsure{\circ|a}}}$] de supposer \uline{en tous cas} $G$ plat sur $S$, \deleted{et $G_1$ \unsure{connexe}} et de plus
% archive sheet number: 31
